\section{Tools}\label{sec:tools}

\subsection{Regular variation}\label{sec:rv}

We use the following facts about a function $T$ that is regularly varying with
index $-\alpha\in(-1,0)$, all from~\cite[Sections~1.2 to~1.5]{bgt1989}.
\begin{enumerate}[(R1)]
\item (Uniform convergence.) $T(\lambda x)/T(x)\to\lambda^{-\alpha}$ uniformly for
$\lambda$ in compact subsets of $(0,\infty)$.
\item (Potter.) For every $c\ge1$ there is $C_c$ with $T(x/c)\le C_cT(x)$ for $x\ge c$.
\item (Karamata.) $\sum_{n\le w}T(n)\sim wT(w)/(1-\alpha)$ and
$\sum_{1\le i\le w}1/(iT(i))\sim1/(\alpha T(w))$ as $w\to\infty$.
\item (Dyadic sums.) $\sum_{l\ge1,\,2^l\le x}2^lT(2^{l-1})\le CxT(x)$ and
$\sum_{l\ge1,\,2^l>x}T(2^{l-1})\le CT(x)$ for $x\ge1$.
\end{enumerate}
(R4) follows from (R2) and Potter's theorem~\cite[Theorem~1.5.6]{bgt1989}, since $wT(w)$ is regularly varying with the
positive index $1-\alpha$ and $T$ with the negative index $-\alpha$, so both sums are
dominated by geometric series (Lemma~\ref{lem:rv-facts}). Under (RV$_\alpha$) the two tails are comparable,
$c\,T_\plus\le T_\minus\le C\,T_\plus$, and $mT_\plus(a_m)\le C$ for all $m\ge1$. We
also use $a_{\lceil m/3\rceil}\ge c\,a_m$, which follows from (R2) for the function
$a$.

\subsection{Cited theorems}\label{sec:cited}

Under (RV$_\alpha$) and (SPAN), $S^\plus_m/a_m$ converges in law to $Z$, and
$S^\minus_m/a_m$ converges to $b^{1/\alpha}Z$, whose density is
$g_b(z)=b^{-1/\alpha}g(zb^{-1/\alpha})$. Note that $t^{-1/\alpha}g_b(yt^{-1/\alpha})=p_{bt}(y)$.

\begin{lemma}[classical inputs]\label{lem:cd}
Assume (RV$_\alpha$) and (SPAN).
\begin{enumerate}[(i)]
\item (Gnedenko's local limit theorem~\cite[Theorem~4.2.1]{ibragimovlinnik1971}.)
$\sup_x|a_mu^\plus_m(x)-g(x/a_m)|\to0$ and $\sup_y|a_mu^\minus_m(y)-g_b(y/a_m)|\to0$ as
$m\to\infty$. Hence $\sup_xu^\sgn_m(x)\le C/a_m$ for all $m\ge1$ and both signs.
\item (Local large deviations~\cite[Theorem~1.1, Eq.~(1.6)]{caravennadoney2019}.)
For both signs there is $C$ with $u^\sgn_m(x)\le C\,\dfrac m{a_m}\,T_\sgn(x)$ for all
$m\ge1$ and $x\ge1$.
\item (The strong renewal theorem.) If $\alpha>1/2$, (SRT$_\sgn$) holds for both signs
\cite{garsialamperti1963,erickson1970}, see~\cite[Theorem~1.4]{caravennadoney2019}.
If $\alpha\le1/2$, (SRT$_\sgn$) holds if and only if
\[
 \lim_{\delta\to0}\limsup_{x\to\infty}\,xT_\sgn(x)\sum_{1\le z\le\delta x}
 \frac{f_\sgn(x-z)}{zT_\sgn(z)^2}=0
\]
\cite[Theorem~1.4]{caravennadoney2019}. In every case (LB$_\sgn$) implies
(SRT$_\sgn$)~\cite{doney1997}, see~\cite[Eq.~(1.15)]{caravennadoney2019}.
\item (The strong renewal theorem as a small-$m$ condition~\cite[Section~4.1,
Eq.~(4.1)]{caravennadoney2019}.) For each sign, (SRT$_\sgn$) holds if and only if
\begin{equation}\label{eq:cd41}
 \lim_{\delta\to0}\limsup_{x\to\infty}I_\sgn(\delta,x)=0,\qquad
 I_\sgn(\delta,x)=xT_\sgn(x)\sum_{1\le m\le\delta/T_\sgn(x)}u^\sgn_m(x).
\end{equation}
\item (Doney's local bound~\cite[Theorem~2.4]{berger2019ptrf}, extending~\cite[Theorem~A]{doney1997}.)
Under (LB$_\sgn$) there is $C$ with $u^\sgn_m(x)\le C\,mT_\sgn(x)/x$ for all $m\ge1$
and $x\ge a_m$.
\end{enumerate}
\end{lemma}

We comment on the form of these statements. In (ii) and (iv) Caravenna and Doney
use the tail $\Prob(X>x)$ and a function $A$ with $A(x)\sim1/\Prob(X>x)$ for renewal
laws. The difference between $\Prob(X>x)$ and $T(x)=\Prob(X\ge x)$ is
$f(x)=o(T(x))$. In (iv) their condition reads
$\lim_{\delta\to0}\limsup_x(x/A(x))\sum_{1\le m\le A(\delta x)}u_m(x)=0$. Since
$A(\delta x)\sim\delta^\alpha/T(x)$ and the sum increases with its upper limit, this
is equivalent to~\eqref{eq:cd41}. (SPAN) makes the span $1$, so their
$\Prob(S_m\in x+(-1,0])$ is $u_m(x)$. In (v) Berger uses a scale $a'_m$ that may differ
from ours by a constant factor. For $c\,a_m\le x\le a_m$ the bound follows from (ii),
since $m/a_m\le m/x$ there, so (v) holds for $x\ge a_m$ with our scale.

\subsection{Elementary bounds}\label{sec:elementary}

The following bounds use only (RV$_\alpha$) and (SPAN). We state them for the plus
sign and they hold for the minus sign with the same proofs. We drop the sign from
the notation.

\begin{lemma}[sums of runs]\label{lem:sums}
\begin{enumerate}[(i)]
\item $\Prob(S_m\ge x)\le C\,mT(x)$ for all $m,x\ge1$.
\item $\Prob(S_m<x)\le e^{-mT(x)}$ for all $m\ge0$ and $x\ge1$.
\item $u_m(x)\le C\,a_m^{-1}e^{-mT(x+1)/3}$ for all $m\ge2$ and $x\ge0$.
\end{enumerate}
\end{lemma}

\begin{proof}
(i) If $S_m\ge x$, then some $X_i\ge x$, or $\sum_iX_i\ones_{\{X_i<x\}}\ge x$. By the
union bound and Markov's inequality,
$\Prob(S_m\ge x)\le mT(x)+m\,\E[X;X<x]/x$. Now
$\E[X;X<x]=\sum_{1\le n<x}\Prob(n\le X<x)\le\sum_{n<x}T(n)\le CxT(x)$ by (R3).

(ii) $S_m<x$ forces $X_i<x$ for every $i\le m$, which has probability
$(1-T(x))^m\le e^{-mT(x)}$.

(iii) Let $h=\lfloor m/2\rfloor\ge m/3$. Conditioning on $S_h$,
$u_m(x)=\sum_{y\le x}u_h(y)u_{m-h}(x-y)\le\Prob(S_h\le x)\sup_zu_{m-h}(z)$. By (ii)
$\Prob(S_h\le x)=\Prob(S_h<x+1)\le e^{-mT(x+1)/3}$, and by Lemma~\ref{lem:cd}(i)
$\sup_zu_{m-h}(z)\le C/a_{m-h}\le C/a_m$, since $m-h\ge m/2$.
\end{proof}

\begin{lemma}[the covering kernel at scale $N$]\label{lem:kernel-limit}
Fix $\delta\in(0,1)$. Uniformly for $t\in[\delta,1/\delta]$ and $x\in[\delta,1]$,
\[
 q^\plus_{m_N(t)+1}(\lfloor xN\rfloor)=T_\plus(N)\bigl[Q_t(x)+o(1)\bigr],\qquad
 q^\minus_{m_N(t)}(\lfloor xN\rfloor)=T_\plus(N)\bigl[Q^b_t(x)+o(1)\bigr].
\]
\end{lemma}

\begin{proof}
Write $m=m_N(t)$ and $k=\lfloor xN\rfloor$, and fix $\theta\in(0,\delta/2)$. Split
$q^\plus_{m+1}(k)=\sum_{r<k}u^\plus_m(r)T_\plus(k-r)$ at $r=k-\theta N$. On
$r\le k-\theta N$ we have $(k-r)/N\in[\theta,1]$, so by (R1)
$T_\plus(k-r)/T_\plus(N)=((k-r)/N)^{-\alpha}+o(1)$ uniformly. By
Lemma~\ref{lem:cd}(i), $u^\plus_m(r)=a_m^{-1}[g(r/a_m)+o(1)]$ uniformly in $r$. The error
terms contribute at most $o(1/a_m)\cdot k\cdot T_\plus(\theta N)=o(T_\plus(N))$, since
$N/a_m$ is bounded. The main term is
$T_\plus(N)\sum_{r\le k-\theta N}a_m^{-1}g(r/a_m)((k-r)/N)^{-\alpha}$, a Riemann sum with mesh
$1/N$ in $z=r/N$. Since $a_m/N\to t^{1/\alpha}$ uniformly and $g$ is bounded and uniformly
continuous, it converges to $\int_0^{x-\theta}p_t(z)(x-z)^{-\alpha}dz$, uniformly in
$t$ and $x$. On $r>k-\theta N$, Lemma~\ref{lem:cd}(i) and (R3) give at most
$(C/a_m)\sum_{i<\theta N+1}T_\plus(i)\le C_\delta\,\theta^{(1-\alpha)/2}T_\plus(N)$ for large
$N$, using (R1) and (R2) to compare $\theta NT_\plus(\theta N)$ with $NT_\plus(N)$. By the
first bound in~\eqref{eq:p-bound} the integral over $[x-\theta,x]$ is at most
$C_\delta\theta^{1-\alpha}$. Let $N\to\infty$ and then $\theta\to0$. The minus kernel is the
same with $u^\minus_{m-1}$ in place of $u^\plus_m$, the density $g_b$ in place of $g$,
and $T_\minus(N)\sim bT_\plus(N)$. The shift from $m$ to $m-1$ does not change the
limit.
\end{proof}

The next bound is the key estimate of the paper. It bounds the covering kernel by
the tail at the level, uniformly in the number of runs, under regular variation
alone.

\begin{lemma}[the key bound]\label{lem:key}
For both signs there is $C$ with $q^\sgn_{m+1}(k)\le C\,T_\sgn(k)$ for all $m\ge0$ and
$k\ge1$.
\end{lemma}

\begin{proof}
We drop the sign. For $m=0$, $q_1=T$. Let $m\ge1$. By~\eqref{eq:q-cover},
$q_{m+1}(k)=\sum_{r<k}u_m(r)T(k-r)$. The part with $r<k/2$ is at most
$T(k/2)\le CT(k)$ by (R2). For $r\ge k/2$ put $i=k-r\in[1,k/2]$, and split the range
of $i$ into dyadic blocks $[w/2,w)$ with $w=2^l$, $l\ge1$, $w\le k$. On a block,
\[
 \sum_{w/2\le i<w}u_m(k-i)T(i)\le T(w/2)\,\Prob\bigl(S_m\in(k-w,k-w/2]\bigr).
\]
The window has at most $w$ lattice points, all at least $k/2$. By
Lemma~\ref{lem:cd}(ii) and (R2) each point has probability at most
$C(m/a_m)T(k)$. By Lemma~\ref{lem:sums}(i) and (R2) the whole window has probability
at most $C\,mT(k)$. So the block is at most
$C\,mT(k)\,T(w/2)\min\{1,w/a_m\}$. Summing with (R4), the blocks with $w\le a_m$ give
at most $C\frac{mT(k)}{a_m}\sum_{w\le a_m}wT(w/2)\le C\,mT(k)T(a_m)$, and the blocks
with $w>a_m$ give at most $C\,mT(k)\sum_{w>a_m}T(w/2)\le C\,mT(k)T(a_m)$. Since
$mT(a_m)\le C$, the total is at most $CT(k)$.
\end{proof}

Note that Lemma~\ref{lem:cd}(ii) is used here at points $x\ge k/2$ that may lie far
below $a_m$ as well as above it, and it holds for all $x$.

\begin{lemma}[decay in $m$]\label{lem:decay}
For both signs there are $C,c$ with $q^\sgn_m(j)\le C\,T_\sgn(j)\,e^{-c\,mT_\sgn(j)}$ for
all $m,j\ge1$.
\end{lemma}

\begin{proof}
We drop the sign. If $mT(j)\le9$, this is Lemma~\ref{lem:key} with $c=1/9$. Let
$mT(j)>9$, so $m\ge10$, and put $h=\lfloor(m-1)/2\rfloor\ge m/3$. Splitting
$S_{m-1}=S_h+(S_{m-1}-S_h)$ gives $q_m(j)=\sum_{s<j}u_h(s)\,q_{m-h}(j-s)$. For $s<j/2$
we use $q_{m-h}(j-s)\le CT(j-s)\le CT(j)$ (Lemma~\ref{lem:key} and (R2)) and
$\Prob(S_h<j/2)\le e^{-hT(j)}\le e^{-mT(j)/3}$ (Lemma~\ref{lem:sums}(ii)). For
$j/2\le s<j$ we use
\[
 \sum_{j/2\le s<j}u_h(s)\,q_{m-h}(j-s)\le\sup_{s<j}u_h(s)\sum_{i\le j/2}CT(i)
 \le C\,a_h^{-1}e^{-hT(j)/3}\,jT(j),
\]
by Lemma~\ref{lem:sums}(iii) and (R3). Since $hT(j)\ge mT(j)/3>3$, we have
$T(j)\ge3/h\ge T(a_h)$ for large $h$, so $j\le Ca_h$ and $jT(j)/a_h\le CT(j)$. Both
parts are at most $CT(j)e^{-mT(j)/9}$.
\end{proof}
